\documentclass[10pt,a4paper]{amsart}
\usepackage{amsmath,amssymb,amsthm,graphicx}
\usepackage{verbatim}
\newenvironment{reference}{\comment}{\endcomment}
%\newenvironment{reference}{}{}
\newenvironment{slogan}{\comment}{\endcomment}
\newenvironment{history}{\comment}{\endcomment}

% For commutative diagrams we use Xy-pic
\usepackage[all]{xy}

% We use 2cell for 2-commutative diagrams.
\xyoption{2cell}
% Theorem environments.
%
\theoremstyle{plain}
\newtheorem{theorem}[subsection]{Theorem}
\newtheorem{proposition}[subsection]{Proposition}
\newtheorem{lemma}[subsection]{Lemma}

\theoremstyle{definition}
\newtheorem{definition}[subsection]{Definition}
\newtheorem{example}[subsection]{Example}
\newtheorem{exercise}[subsection]{Exercise}
\newtheorem{situation}[subsection]{Situation}

\theoremstyle{remark}
\newtheorem{remark}[subsection]{Remark}
\newtheorem{remarks}[subsection]{Remarks}

\numberwithin{equation}{subsection}

% Macros
%
\def\lim{\mathop{\mathrm{lim}}\nolimits}
\def\colim{\mathop{\mathrm{colim}}\nolimits}
\def\Spec{\mathop{\mathrm{Spec}}}
\def\Hom{\mathop{\mathrm{Hom}}\nolimits}
\def\Ext{\mathop{\mathrm{Ext}}\nolimits}
\def\SheafHom{\mathop{\mathcal{H}\!\mathit{om}}\nolimits}
\def\SheafExt{\mathop{\mathcal{E}\!\mathit{xt}}\nolimits}
\def\Sch{\mathit{Sch}}
\def\Mor{\mathop{\mathrm{Mor}}\nolimits}
\def\Ob{\mathop{\mathrm{Ob}}\nolimits}
\def\Sh{\mathop{\mathit{Sh}}\nolimits}
\def\NL{\mathop{N\!L}\nolimits}
\def\CH{\mathop{\mathrm{CH}}\nolimits}
\def\proetale{{pro\text{-}\acute{e}tale}}
\def\etale{{\acute{e}tale}}
\def\QCoh{\mathit{QCoh}}
\def\Ker{\mathop{\mathrm{Ker}}}
\def\Im{\mathop{\mathrm{Im}}}
\def\Coker{\mathop{\mathrm{Coker}}}
\def\Coim{\mathop{\mathrm{Coim}}}

% Boxtimes
%
\DeclareMathSymbol{\boxtimes}{\mathbin}{AMSa}{"02}

%
% Macros for moduli stacks/spaces
%
\def\QCohstack{\mathcal{QC}\!\mathit{oh}}
\def\Cohstack{\mathcal{C}\!\mathit{oh}}
\def\Spacesstack{\mathcal{S}\!\mathit{paces}}
\def\Quotfunctor{\mathrm{Quot}}
\def\Hilbfunctor{\mathrm{Hilb}}
\def\Curvesstack{\mathcal{C}\!\mathit{urves}}
\def\Polarizedstack{\mathcal{P}\!\mathit{olarized}}
\def\Complexesstack{\mathcal{C}\!\mathit{omplexes}}
% \Pic is the operator that assigns to X its picard group, usage \Pic(X)
% \Picardstack_{X/B} denotes the Picard stack of X over B
% \Picardfunctor_{X/B} denotes the Picard functor of X over B
\def\Pic{\mathop{\mathrm{Pic}}\nolimits}
\def\Picardstack{\mathcal{P}\!\mathit{ic}}
\def\Picardfunctor{\mathrm{Pic}}
\def\Deformationcategory{\mathcal{D}\!\mathit{ef}}
\usepackage{cleveref}
\crefname{originallemma}{original lemma}{original lemmas}
\Crefname{originallemma}{Original lemma}{Original lemmas}
\makeatletter
\newlabel{coherent-lemma-coherent-proper-ample}{{0B5T}{source}{}{sourcepointer-1}{}}
\newlabel{coherent-lemma-coherent-proper-ample@cref}{{[originallemma][0][]0B5T}{[1][0][]source}{}{}{}}
\makeatother
\crefname{originalreference}{original reference}{original references}
\Crefname{originalreference}{Original reference}{Original references}
\makeatletter
\newlabel{definition-genus}{{0B9W}{source}{}{sourcepointer-2}{}}
\newlabel{definition-genus@cref}{{[originalreference][0][]0B9W}{[1][0][]source}{}{}{}}
\makeatother
\crefname{originallemma}{original lemma}{original lemmas}
\Crefname{originallemma}{Original lemma}{Original lemmas}
\makeatletter
\newlabel{divisors-lemma-characterize-OD}{{01X0}{source}{}{sourcepointer-3}{}}
\newlabel{divisors-lemma-characterize-OD@cref}{{[originallemma][0][]01X0}{[1][0][]source}{}{}{}}
\makeatother
\crefname{originallemma}{original lemma}{original lemmas}
\Crefname{originallemma}{Original lemma}{Original lemmas}
\makeatletter
\newlabel{divisors-lemma-finite-trivialize-invertible-upstairs}{{0BUT}{source}{}{sourcepointer-4}{}}
\newlabel{divisors-lemma-finite-trivialize-invertible-upstairs@cref}{{[originallemma][0][]0BUT}{[1][0][]source}{}{}{}}
\makeatother
\crefname{originalremark}{original remark}{original remarks}
\Crefname{originalremark}{Original remark}{Original remarks}
\makeatletter
\newlabel{divisors-remark-ses-regular-section}{{0C6K}{source}{}{sourcepointer-5}{}}
\newlabel{divisors-remark-ses-regular-section@cref}{{[originalremark][0][]0C6K}{[1][0][]source}{}{}{}}
\makeatother
\crefname{originallemma}{original lemma}{original lemmas}
\Crefname{originallemma}{Original lemma}{Original lemmas}
\makeatletter
\newlabel{varieties-lemma-ample-curve}{{0B5X}{source}{}{sourcepointer-6}{}}
\newlabel{varieties-lemma-ample-curve@cref}{{[originallemma][0][]0B5X}{[1][0][]source}{}{}{}}
\makeatother
\crefname{originallemma}{original lemma}{original lemmas}
\Crefname{originallemma}{Original lemma}{Original lemmas}
\makeatletter
\newlabel{varieties-lemma-degree-effective-Cartier-divisor}{{0AYY}{source}{}{sourcepointer-7}{}}
\newlabel{varieties-lemma-degree-effective-Cartier-divisor@cref}{{[originallemma][0][]0AYY}{[1][0][]source}{}{}{}}
\makeatother
\title{A g\textasciicircum{}r\_d on a proper curve with H\textasciicircum{}0(O\_X) = k satisfies r <= d, with equality forcing genus 0}
\author{The Stacks Project Authors}
\date{stacks-project:curves.tex:local-snapshot-20261012; GNU FDL 1.2 or later}
\begin{document}
\maketitle
\noindent\textit{Editorial note.} All delivered mathematical text is the author's own, copied from the pinned Stacks Project chapter file curves.tex; no statement or proof step was rewritten, reordered or removed.\par\medskip
\noindent\textit{Editorial note.} The counted proof is the author's complete argument: the section-sequence dimension count with the triviality of L|\_D, the equality-case strategy remark, the doubling step with its explicit dimension identity, and the closing ampleness argument are all retained verbatim, including the author's footnote on the triviality of L|\_D.\par\medskip
\noindent\textit{Editorial note.} The label lemma-grd-inequalities names the counted statement, and the proof environment immediately following it in the source is delivered verbatim as its proof.\par\medskip
\noindent\textit{Editorial note.} The counted statement is carried by the author's own lemma environment in the chapter source.\par\medskip
\noindent\textit{Allowed standard prerequisites (editorial).}\begin{itemize}
\item Cartier divisors from sections. Divisors, Lemma divisors-lemma-characterize-OD and Remark divisors-remark-ses-regular-section; Varieties, Lemma varieties-lemma-degree-effective-Cartier-divisor. Supplies the sequence 0 -> O\_X -> L -> L|\_D -> 0 for the zero divisor D of a section, with deg(D) = deg(L) = d.
\item Ampleness and vanishing on curves. Varieties, Lemma varieties-lemma-ample-curve; Cohomology of Schemes, Lemma coherent-lemma-coherent-proper-ample. Provides the power L\textasciicircum{}tensor n with vanishing H\textasciicircum{}1 used to close the equality case.
\end{itemize}\medskip
\ifdefined\hypertarget\hypertarget{sourcepointer-1}{}\fi
\noindent\textit{Source reference.} Original lemma 0B5T: sources/stacks-chow/tags.txt, source lines 14046--14046. Prerequisite pointer invoked at proof line(s) 731; it names a standard fact of a sibling chapter that the counted argument uses, and every step of the counted proof built on it is retained verbatim in the delivered body.\par\medskip
\ifdefined\hypertarget\hypertarget{sourcepointer-2}{}\fi
\noindent\textit{Source reference.} Original reference 0B9W: sources/stacks-curves/pinned-original-source.tex, source lines 1677--1677. Prerequisite pointer invoked at proof line(s) 664; it names a standard fact of the same chapter that the counted argument uses, and every step of the counted proof built on it is retained verbatim in the delivered body.\par\medskip
\ifdefined\hypertarget\hypertarget{sourcepointer-3}{}\fi
\noindent\textit{Source reference.} Original lemma 01X0: sources/stacks-chow/tags.txt, source lines 2501--2501. Prerequisite pointer invoked at proof line(s) 676; it names a standard fact of a sibling chapter that the counted argument uses, and every step of the counted proof built on it is retained verbatim in the delivered body.\par\medskip
\ifdefined\hypertarget\hypertarget{sourcepointer-4}{}\fi
\noindent\textit{Source reference.} Original lemma 0BUT: sources/stacks-chow/tags.txt, source lines 14886--14886. Prerequisite pointer invoked at proof line(s) 683; it names a standard fact of a sibling chapter that the counted argument uses, and every step of the counted proof built on it is retained verbatim in the delivered body.\par\medskip
\ifdefined\hypertarget\hypertarget{sourcepointer-5}{}\fi
\noindent\textit{Source reference.} Original remark 0C6K: sources/stacks-chow/tags.txt, source lines 15300--15300. Prerequisite pointer invoked at proof line(s) 677; it names a standard fact of a sibling chapter that the counted argument uses, and every step of the counted proof built on it is retained verbatim in the delivered body.\par\medskip
\ifdefined\hypertarget\hypertarget{sourcepointer-6}{}\fi
\noindent\textit{Source reference.} Original lemma 0B5X: sources/stacks-chow/tags.txt, source lines 14050--14050. Prerequisite pointer invoked at proof line(s) 728; it names a standard fact of a sibling chapter that the counted argument uses, and every step of the counted proof built on it is retained verbatim in the delivered body.\par\medskip
\ifdefined\hypertarget\hypertarget{sourcepointer-7}{}\fi
\noindent\textit{Source reference.} Original lemma 0AYY: sources/stacks-chow/tags.txt, source lines 13806--13806. Prerequisite pointer invoked at proof line(s) 678; it names a standard fact of a sibling chapter that the counted argument uses, and every step of the counted proof built on it is retained verbatim in the delivered body.\par\medskip
\begin{lemma}
\label{lemma-grd-inequalities}
Let $k$ be a field. Let $X$ be a proper curve over $k$ with
$H^0(X, \mathcal{O}_X) = k$. If $X$ has a $\mathfrak g^r_d$, then
$r \leq d$. If equality holds, then $H^1(X, \mathcal{O}_X) = 0$, i.e.,
the genus of $X$ (Definition \ref{definition-genus}) is $0$.
\end{lemma}

\begin{proof}
Let $(\mathcal{L}, V)$ be a $\mathfrak g^r_d$. Since this will only
increase $r$, we may assume $V = H^0(X, \mathcal{L})$. Choose a
nonzero element $s \in V$.  Then the zero scheme of $s$ is an effective Cartier
divisor $D \subset X$, we have $\mathcal{L} = \mathcal{O}_X(D)$, and
we have a short exact sequence
$$
0 \to \mathcal{O}_X \to \mathcal{L} \to \mathcal{L}|_D \to 0
$$
see Divisors, Lemma \ref{divisors-lemma-characterize-OD} and
Remark \ref{divisors-remark-ses-regular-section}.
By Varieties, Lemma \ref{varieties-lemma-degree-effective-Cartier-divisor}
we have $\deg(D) = \deg(\mathcal{L}) = d$.
Since $D$ is an Artinian scheme we have
$\mathcal{L}|_D \cong \mathcal{O}_D$\footnote{In our case this
follows from Divisors, Lemma
\ref{divisors-lemma-finite-trivialize-invertible-upstairs}
as $D \to \Spec(k)$ is finite.}.
Thus
$$
\dim_k H^0(D, \mathcal{L}|_D) = \dim_k H^0(D, \mathcal{O}_D) = \deg(D) = d
$$
On the other hand, by assumption
$\dim_k H^0(X, \mathcal{O}_X) = 1$ and $\dim H^0(X, \mathcal{L}) = r + 1$.
We conclude that $r + 1 \leq 1 + d$, i.e., $r \leq d$ as in the lemma.

\medskip\noindent
Assume equality holds. Then
$H^0(X, \mathcal{L}) \to H^0(X, \mathcal{L}|_D)$ is surjective.
If we knew that $H^1(X, \mathcal{L})$ was zero, then we would
conclude that $H^1(X, \mathcal{O}_X)$ is zero by the long exact
cohomology sequence and the proof would
be complete. Our strategy will be to replace $\mathcal{L}$ by a
large power which has vanishing. As $\mathcal{L}|_D$ is the
trivial invertible module (see above), we can
find a section $t$ of $\mathcal{L}$ whose restriction
of $D$ generates $\mathcal{L}|_D$.
Consider the multiplication map
$$
\mu :
H^0(X, \mathcal{L}) \otimes_k H^0(X, \mathcal{L})
\longrightarrow
H^0(X, \mathcal{L}^{\otimes 2})
$$
and consider the short exact sequence
$$
0 \to \mathcal{L} \xrightarrow{s}
\mathcal{L}^{\otimes 2} \to \mathcal{L}^{\otimes 2}|_D \to 0
$$
Since $H^0(\mathcal{L}) \to H^0(\mathcal{L}|_D)$ is surjective and since
$t$ maps to a trivialization of $\mathcal{L}|_D$ we see that
$\mu(H^0(X, \mathcal{L}) \otimes t)$ gives a subspace of
$H^0(X, \mathcal{L}^{\otimes 2})$ surjecting onto the global sections of
$\mathcal{L}^{\otimes 2}|_D$. Thus we see that
$$
\dim H^0(X, \mathcal{L}^{\otimes 2}) = r + 1 + d = 2r + 1 =
\deg(\mathcal{L}^{\otimes 2}) + 1
$$
Ok, so $\mathcal{L}^{\otimes 2}$ has the same property as $\mathcal{L}$, i.e.,
that the dimension of the space of global sections is equal to the
degree plus one. Since $\mathcal{L}$ is ample
(Varieties, Lemma \ref{varieties-lemma-ample-curve})
there exists some $n_0$ such that $\mathcal{L}^{\otimes n}$
has vanishing $H^1$ for all $n \geq n_0$
(Cohomology of Schemes, Lemma \ref{coherent-lemma-coherent-proper-ample}).
Thus applying the argument above to $\mathcal{L}^{\otimes n}$
with $n = 2^m$ for some sufficiently large $m$ we conclude the
lemma is true.
\end{proof}

\end{document}
